\documentclass[10pt]{article}
\usepackage[T1]{fontenc}
\usepackage{lmodern}
\usepackage[a4paper,margin=0.86in]{geometry}
\usepackage{amsmath,amssymb,amsthm,mathtools,amscd,enumitem,xcolor,hyperref}
\hypersetup{colorlinks=true,linkcolor=blue!55!black,citecolor=blue!55!black,urlcolor=blue!55!black}
\setlength{\emergencystretch}{3em}
\setlength{\parindent}{0pt}
\setlength{\parskip}{0.52em}
\setlist{leftmargin=1.65em,itemsep=0.24em,topsep=0.25em}
\sloppy

\newtheorem{theorem}{Theorem}[section]
\newtheorem{proposition}[theorem]{Proposition}
\newtheorem{lemma}[theorem]{Lemma}
\newtheorem{corollary}[theorem]{Corollary}
\theoremstyle{definition}
\newtheorem{definition}[theorem]{Definition}
\newtheorem{convention}[theorem]{Convention}
\newtheorem{remark}[theorem]{Remark}

\newcommand{\Fq}{\mathbf F_q}
\newcommand{\Fbar}{\overline{\mathbf F}_q}
\newcommand{\Ql}{\mathbf Q_\ell}
\newcommand{\Qlb}{\overline{\mathbf Q}_\ell}
\newcommand{\Dbc}{D^b_c}
\newcommand{\DW}{D^b_{\mathrm W,c}}
\newcommand{\WtLe}{\operatorname{Wt}_{\leq}}
\newcommand{\WtGe}{\operatorname{Wt}_{\geq}}
\newcommand{\IntWtLe}{\operatorname{Wt}^{\mathbf Z}_{\leq}}
\newcommand{\IntWtGe}{\operatorname{Wt}^{\mathbf Z}_{\geq}}
\newcommand{\Fr}{\operatorname{Fr}}
\newcommand{\Supp}{\operatorname{Supp}}
\newcommand{\Spec}{\operatorname{Spec}}
\newcommand{\dep}[1]{\par\smallskip\noindent\textit{Direct dependencies.} #1\par}

\begin{document}
\begin{center}
{\large\bfseries Deligne's Global Direct-Image Weight Theorem\par}
\medskip
{\normalsize Devissage through curves and the affine line\par}
\end{center}

\begin{abstract}
Let \(f:X_0\to Y_0\) be separated and of finite type over a finite field.
Deligne's theorem asserts that \(Rf_!\) preserves upper weights, with the
cohomological shift built into the definition of the weight of a complex.
This note records a theorem-by-theorem route through finite stratification,
localization, finite and radicial descent, the Leray spectral sequence,
relative curves, and the affine-line estimate.  The unconditional route is
stated for mixed complexes whose pure graded pieces have integral weights.
The separate real-weight statement retains the additional local-Fourier
hypotheses under which it is valid.  All categories below are ordinary
triangulated categories.
\end{abstract}

\tableofcontents

\section{Coefficient categories and Frobenius}

\subsection{Constructible pro-etale coefficients}

Fix \(k=\Fq\) of characteristic \(p\), a prime \(\ell\ne p\), an algebraic
closure \(\bar k=\Fbar\), and an embedding
\[
                       \iota:\Qlb\hookrightarrow\mathbf C.
\]
Fix a nontrivial additive character \(\psi:k\to\Qlb^\times\).
Every scheme in the note is of finite type over \(k\).  If \(X_0\) is such
a scheme, write \(X=X_0\times_k\bar k\).

Let \(E/\Ql\) be finite, let \(\mathcal O_E\) be its valuation ring, and
choose a uniformizer \(\varpi\).  On \(X_{\mathrm{pro\acute et}}\), put
\[
 E_X=\left(\varprojlim_n(\mathcal O_E/\varpi^n)_X\right)[1/\varpi].
\]
Put \(A_n=\mathcal O_E/\varpi^n\).  An \((A_n)_X\)-module sheaf is
finite-torsion constructible if it is locally constant with finite,
equivalently finitely generated, \(A_n\)-module fibers on every member of a
finite locally closed constructible partition.  Denote the resulting
abelian category by \(\mathsf{Cons}(X,A_n)\), and define
\[
 D^b_c(X,A_n)=
 \{K\in D(X_{\mathrm{pro\acute et}},(A_n)_X):
   \mathcal H^i(K)\in\mathsf{Cons}(X,A_n),\ 
   \mathcal H^i(K)=0\text{ for }|i|\gg0\}.
\]
A sheaf of \(E_X\)-modules is \emph{lisse of finite rank} if it is
pro-etale locally isomorphic to \(E_X^r\) for some finite \(r\).  It is
\emph{constructible} if \(X\) has a finite locally closed stratification on
which the restrictions are lisse of finite rank.  Define
\[
 \Dbc(X,E)=
 \{K\in D(X_{\mathrm{pro\acute et}},E_X):
 K\text{ is bounded and every }\mathcal H^m(K)\text{ is constructible}\}.
\]
It is the full ordinary triangulated subcategory of the unbounded derived
category with the displayed property.

\begin{convention}[Fourier coefficient field]
Whenever \(\mathcal L_\psi\) or \(\operatorname{FT}_\psi\) occurs, replace
\(E\) by a finite extension containing the finite set \(\psi(k)\); the
coefficient field is then fixed for that statement.
\end{convention}

\begin{definition}[Geometric Frobenius]
The arithmetic field Frobenius is
\(\Fr_{\mathrm{ar}}:a\mapsto a^q\), and the geometric Galois Frobenius is
\(\gamma=\Fr_{\mathrm{ar}}^{-1}\).  These are semilinear automorphisms of
the base field.  Separately, the relative \(q\)-power Frobenius followed by
the descent identification of the Frobenius twist defines the scheme
universal homeomorphism
\[
                       \Phi_X:X\longrightarrow X.
\]
Topological invariance and descent give a natural comparison between
\(\Phi_X^*\) on pro-etale sheaves and pullback by the semilinear geometric
Galois action; they do not identify the two morphisms of geometric objects.
For \(\varphi:\Phi_X^*\mathcal F\xrightarrow{\sim}\mathcal F\), put
\[
 \varphi^{[0]}=\operatorname{id},\qquad
 \varphi^{[n+1]}=\varphi^{[n]}\circ(\Phi_X^{[n]})^*\varphi,
 \qquad \Phi_X^{[n]}=\Phi_X^n.
\]
For a closed point \(x\in|X_0|\), set
\(d(x)=[\kappa(x):k]\) and \(q_x=q^{d(x)}\).  The scheme-Frobenius orbit
of a geometric point \(\bar x\to x\) has length \(d(x)\).  The stalk of
\(\varphi^{[d(x)]}\), composed with the canonical point-orbit pullback
identification, is
\[
                  F_x:\mathcal F_{\bar x}\longrightarrow
                       \mathcal F_{\bar x}.
\]
Under the descent comparison this is the local geometric Galois Frobenius.
For \(y\in X_0(\mathbf F_{q^n})\) lying over \(x\), the fixed-point
operator from \(\varphi^{[n]}\) is conjugate to \(F_x^{n/d(x)}\).
\end{definition}

\begin{definition}[Constructible Weil sheaves]
A constructible \(E\)-Weil sheaf on \(X_0\) is a pair
\((\mathcal F,\varphi_{\mathcal F})\), where \(\mathcal F\) is a
constructible \(E_X\)-sheaf and
\[
             \varphi_{\mathcal F}:\Phi_X^*\mathcal F
                         \xrightarrow{\sim}\mathcal F.
\]
Morphisms commute with these isomorphisms.  The resulting abelian category
is denoted \(\mathsf{Cons}_{\mathrm W}(X_0,E)\).
\end{definition}

\begin{definition}[Strict Weil derived category]
The strict Weil derived category is the ordinary Verdier localization
\[
 \DW(X_0,E)=
 K^b(\mathsf{Cons}_{\mathrm W}(X_0,E))
 [\text{equivariant quasi-isomorphisms}^{-1}].
\]
Its distinguished triangles are the images of mapping-cone triangles in
the localization.  The exact forgetful functor
\[
       \operatorname{forg}_X:\DW(X_0,E)\longrightarrow\Dbc(X,E)
\]
forgets Frobenius.  In particular, a Weil structure is attached before
localizing; it is not merely an isomorphism attached to an object of the
already formed derived category.
\end{definition}

\subsection{Exceptional and ordinary direct image}

Let \(f_0:X_0\to Y_0\) be separated and of finite type and let
\(f:X\to Y\) be its geometric base change.  A compactification is a
factorization
\[
                 X\xrightarrow{j}\overline X\xrightarrow{\bar f}Y
\]
with \(j\) open and \(\bar f\) proper.  On the unbounded derived category
define
\[
                Rf_!^{\mathrm{amb}}=R\bar f_*\,j_!.
\]
The compactification comparison theorem gives canonical isomorphisms
between different choices, satisfying the cocycle condition.  Thus the
notation \(Rf_!^{\mathrm{amb}}\) is independent of the selected
compactification.

\begin{theorem}[Bounded constructibility of direct images]
The ambient functors restrict to exact functors of ordinary triangulated
categories
\[
 Rf_!,Rf_*:\Dbc(X,E)\longrightarrow\Dbc(Y,E).
\]
For composable separated finite-type maps \(X\xrightarrow fY\xrightarrow gZ\),
there are coherent natural isomorphisms
\[
 R(gf)_!\xrightarrow{\sim}Rg_!Rf_!,\qquad
 R(gf)_*\xrightarrow{\sim}Rg_*Rf_*.
\]
The restrictions commute with finite extension of \(E\), with shifts and
distinguished triangles, and with the usual base-change morphisms.
\end{theorem}

\begin{theorem}[Strict Frobenius lift]
There are exact functors
\[
 Rf_{!,\mathrm W},Rf_{*,\mathrm W}:\DW(X_0,E)
                       \longrightarrow\DW(Y_0,E)
\]
whose underlying functors are \(Rf_!\) and \(Rf_*\).  The Frobenius on the
first is the composite
\[
 \Phi_Y^*Rf_!K\xrightarrow{\sim}Rf_!\Phi_X^*K
             \xrightarrow{Rf_!(\varphi_K)}Rf_!K,
\]
where the first arrow is compact-support base change followed by topological
invariance under the relative Frobenius.  The corresponding formula with
\(Rf_*\) defines the second lift.  These lifts are compatible with
composition, compactification comparison, and scalar extension.
\end{theorem}

For \(K\in\DW(X_0,E)\), write
\[
              R^mf_!K=\mathcal H^m(Rf_{!,\mathrm W}K),\qquad
              R^mf_*K=\mathcal H^m(Rf_{*,\mathrm W}K).
\]
These are constructible Weil sheaves by the preceding two theorems.

\begin{definition}[Derived internal Hom]
For \(K,L\in D(X_{\mathrm{pro\acute et}},E_X)\), choose a K-flat
quasi-isomorphism \(P_K\to K\) and a K-injective quasi-isomorphism
\(L\to I_L\).  Thus \(P_K\otimes A\) is acyclic and
\(\mathcal Hom(A,I_L)\) is acyclic for every acyclic complex \(A\).
Define
\[
 R\mathcal Hom_E(K,L)=\mathcal Hom_E^\bullet(P_K,I_L),
\]
where
\[
 \mathcal Hom_E^n(P_K,I_L)=\prod_p
       \mathcal Hom_E(P_K^p,I_L^{p+n}),
 \qquad d(f)=d_{I_L}f-(-1)^nf d_{P_K}.
\]
This construction is independent of the replacements in the ambient
ordinary derived category.
\end{definition}

\begin{definition}[Ambient equivariant derived category]
Let \(\mathsf{Mod}_{\mathrm W}(X_0,E)\) be the abelian category of all
\(E_X\)-module sheaves \(M\), without a constructibility hypothesis,
equipped with \(\Phi_X^*M\xrightarrow{\sim}M\).  Let
\(D_{\mathrm W}(X_0,E)=D(\mathsf{Mod}_{\mathrm W}(X_0,E))\), and let
\(D^b_{\mathrm W,c,\mathrm{amb}}(X_0,E)\) be its full triangulated
subcategory of objects with bounded constructible underlying cohomology.
Inclusion induces
\[
 \operatorname{Real}_{\mathrm W}:\DW(X_0,E)
             \longrightarrow D^b_{\mathrm W,c,\mathrm{amb}}(X_0,E).
\]
\end{definition}

\begin{theorem}[Bounded strict equivariant realization]
The functor \(\operatorname{Real}_{\mathrm W}\) is an equivalence of
ordinary triangulated categories, compatible with forgetting Frobenius,
finite scalar extension, shifts, cones, and cohomology sheaves.
\dep{The Serre property of constructible equivariant sheaves, equivariant
resolutions in \(\mathsf{Mod}_{\mathrm W}(X_0,E)\), finite constructible
stratification, and bounded-cohomology truncation.}
\end{theorem}

\begin{theorem}[Bounded constructibility and strict internal Hom]
Derived internal Hom restricts to an exact bifunctor
\[
 R\mathcal Hom_E(-,-):\Dbc(X,E)^{\mathrm{op}}\times\Dbc(X,E)
                                      \longrightarrow\Dbc(X,E).
\]
For strict Weil complexes \((K,\varphi_K)\), \((L,\varphi_L)\), bounded
strict constructible realization gives an exact bifunctor on \(\DW\) whose
Frobenius arrow is the typed composite
\[
 \Phi_X^*R\mathcal Hom_E(K,L)
 \xrightarrow{\sim}R\mathcal Hom_E(\Phi_X^*K,\Phi_X^*L)
 \xrightarrow{R\mathcal Hom_E(\varphi_K^{-1},\varphi_L)}
 R\mathcal Hom_E(K,L).
\]
It commutes with finite scalar extension.
\dep{Equivariant K-flat and K-injective resolutions, finite Tor amplitude,
bounded constructibility of internal Hom, and bounded strict constructible
realization in the ambient equivariant derived category.}
\end{theorem}

\begin{definition}[Tate objects and twists]
On the pro-etale site put
\[
 \mathbf Z_\ell(1)=\varprojlim_n\mu_{\ell^n},
 \qquad
 \mu_{\ell^n}=\ker([\ell^n]:\mathbf G_m\to\mathbf G_m),
\]
where the transition map is \(\zeta\mapsto\zeta^\ell\).  Define
\[
 E(1)=\mathbf Z_\ell(1)\otimes_{\mathbf Z_\ell}E,
 \quad E(m)=E(1)^{\otimes m}\ (m\ge0),
 \quad E(-m)=R\mathcal Hom_E(E(m),E),
\]
 For \(K\in\DW(X_0,E)\), define
\(K(m)=K\otimes_E^LE(m)\).  Functorial pullback of roots of unity and
the descent comparison define compatible arrows
\[
 \Phi_X^*\mu_{\ell^n}\xrightarrow{\sim}\mu_{\ell^n}.
\]
Their inverse limit and rationalization give the strict Weil arrow
\[
                  \varphi_{E(1)}:\Phi_X^*E(1)\xrightarrow{\sim}E(1),
\]
and tensor products and strict duals give the arrows on every \(E(m)\) and
\(K(m)\).
\end{definition}

\begin{definition}[Finite \(Q\)-Frobenius modules]
Let \(Q>1\) be an integral power of \(q\).  A finite \(Q\)-Frobenius module
over \(E\) is a triple
\[
                         (V,F,Q),
\]
where \(V\) is a finite-dimensional \(E\)-vector space and
\(F\in\operatorname{Aut}_E(V)\).  A morphism
\((V,F,Q)\to(V',F',Q)\) is an \(E\)-linear map commuting with Frobenius.
Submodules, quotient modules, filtrations, and graded modules always mean
\(F\)-stable ones.
\end{definition}

\begin{definition}[Weights of a finite \(Q\)-Frobenius module]
Fix \(w\in\mathbf R\).  A finite \(Q\)-Frobenius module \((V,F,Q)\) is
\emph{\(\iota\)-pure of weight \(w\)} if every eigenvalue \(\alpha\) of
\(F\) in \(\Qlb\) satisfies
\[
                           |\iota(\alpha)|=Q^{w/2}.
\]
It is \emph{mixed of \(\iota\)-weights at most \(w\)} if, after one finite
extension of \(E\), it has a finite Frobenius-stable filtration whose
nonzero graded modules are \(\iota\)-pure of real weights at most \(w\).
It is \emph{integrally mixed of \(\iota\)-weights at most \(w\)} if those
graded weights may all be chosen in \(\mathbf Z\).  The corresponding
predicates with ``at least'' use the opposite inequalities.
\end{definition}

\begin{definition}[Fixed-embedding and all-conjugates purity]
Fix \(w\in\mathbf R\).
The purity predicate in the preceding definition is a
\emph{fixed-\(\iota\) predicate}: it constrains only the absolute value
obtained from the chosen embedding \(\iota:\Qlb\hookrightarrow\mathbf C\).
A finite \(Q\)-Frobenius module is \emph{pure of weight \(w\) at all
conjugates} if every Frobenius eigenvalue is algebraic over \(\mathbf Q\)
and, for every embedding of its number field into \(\mathbf C\), the image
of that eigenvalue has absolute value \(Q^{w/2}\).  All-conjugates purity
implies fixed-\(\iota\) purity after identifying the algebraic eigenvalues
inside \(\Qlb\); the converse is not part of the definition.
\end{definition}

\begin{definition}[Twists, duals, and power change]
For a finite \(Q\)-Frobenius module \(M=(V,F,Q)\) and \(m\in\mathbf Z\), put
\[
 M(m)=(V,Q^{-m}F,Q),\qquad
 M^\vee=(V^\vee,(F^{-1})^\vee,Q).
\]
For an integer \(r\geq1\), its \emph{power change} is
\[
                         M^{[r]}=(V,F^r,Q^r).
\]
Thus \(M(m)\) shifts every weight by \(-2m\), \(M^\vee\) negates every
weight, and \(M^{[r]}\) has exactly the same weights as \(M\).  These
identities hold for purity, mixed upper and lower bounds, and integral
mixedness.  A Frobenius-equivariant arrow \(N:M\to M(-1)\) is equivalently
an \(E\)-linear map satisfying \(FNF^{-1}=Q^{-1}N\).
\end{definition}

\begin{proposition}[Tate Frobenius and weight shift]
Let \(m\in\mathbf Z\).  At a closed point \(x\), the ordered
\([\kappa(x):k]\)-fold iterate of \(\varphi_{E(1)}\) acts by
\(q_x^{-1}\).  Hence the Tate twist of the finite \(q_x\)-Frobenius module
\((V,F_x,q_x)\) is
\[
                   (V,F_x,q_x)(m)=(V,q_x^{-m}F_x,q_x),
\]
and it changes every weight by \(-2m\).  In particular,
\((V,F_x,q_x)(-1)\) has Frobenius \(q_xF_x\).
\end{proposition}

\begin{proposition}[Closed-point Frobenius module]
Let \((\mathcal F,\varphi_{\mathcal F})\) be a constructible Weil sheaf and
let \(x\in|X_0|\).  Its closed-point Frobenius module is
\[
                       (\mathcal F_{\bar x},F_x,q_x).
\]
For \(n\geq1\) and \(y\in X_0(\mathbf F_{q^n})\) lying over \(x\), the module defined by
the fixed-point operator at \(y\) is the power change by
\(r=n/d(x)\).  Tate twists, duals, purity, and mixed upper and lower bounds
agree with the preceding definitions under this identification.
\end{proposition}

\subsection{Conditional exceptional duality}

\begin{definition}[Finite-level exceptional-duality hypotheses]
For every separated finite-type \(k\)-scheme \(Z_0\), let
\(a_Z:Z\to\Spec\bar k\) be the geometric structural map.  The finite-level exceptional-duality
hypotheses assert the following.
\begin{enumerate}[label=\textup{(\roman*)}]
\item The ambient functor \(Ra_{Z,!}\) preserves small coproducts on a
derived category satisfying Brown representability, and hence has an exact
right adjoint \(Ra_Z^!\).
\item The finite-level right adjoints \(Ra_{Z,n}^!\) exist, the complexes
\(\omega_{Z,n}=Ra_{Z,n}^!A_n\) belong to \(D^b_c(Z,A_n)\) in one
cohomological interval independent of \(n\), and
\[
 \omega_{Z,n+1}\otimes^L_{A_{n+1}}A_n
                    \xrightarrow{\sim}\omega_{Z,n}.
\]
\item Internal Hom into \(\omega_{Z,n}\) preserves bounded
constructibility and its finite-level evaluation map is an isomorphism.
These objects and maps commute with reduction, derived inverse limit,
rationalization, and finite scalar extension; the normalized systems have
uniformly finite Tor amplitude.
\item Exceptional Frobenius base change and the tensor--Hom, projection,
and trace transformations are compatible at every finite level and with
adic realization.
\end{enumerate}
No rational biduality or direct-image exchange conclusion is included in
these hypotheses.
\end{definition}

\begin{theorem}[Conditional rational exceptional duality]
Let \(Z_0\) be a separated finite-type \(k\)-scheme, put
\(Z=Z_0\times_k\bar k\), and let
\(a_Z:Z\to\Spec\bar k\) be its structural morphism.  Assume the
finite-level exceptional-duality hypotheses for \(Z_0\).  Put
\[
             \omega_Z=Ra_Z^!E.
\]
For \(K\in\Dbc(Z,E)\), define
\[
                   D_ZK=R\mathcal Hom_E(K,\omega_Z).
\]
Then
\(\omega_Z\in\Dbc(Z,E)\), the functor
\(D_Z:\Dbc(Z,E)^{\mathrm{op}}\to\Dbc(Z,E)\) is exact, and the canonical
evaluation map \(K\to D_ZD_ZK\) is an isomorphism.  For every separated
 finite-type \(f_0:Z_0\to T_0\) for which the finite-level
exceptional-duality hypotheses also hold on \(T_0\), put
\(T=T_0\times_k\bar k\) and let \(f:Z\to T\) be the geometric map.  There
are natural isomorphisms
\[
 D_TRf_*\xrightarrow{\sim}Rf_!D_Z,
 \qquad
 D_TRf_!\xrightarrow{\sim}Rf_*D_Z.
\]
They satisfy identity, composition, biduality, and finite-scalar-extension
compatibilities.
\dep{Finite-level constructible Verdier duality, uniform cohomological and
Tor amplitude, adic realization, rationalization, tensor--Hom adjunction,
the projection formula, and the adjunction trace.}
\end{theorem}

\begin{theorem}[Exceptional Frobenius exchange]
Retain \(Z_0,Z,a_Z\), and \(\omega_Z\) from the preceding theorem.  The
adjoint mate of Frobenius base change and
the strict coefficient-unit arrow
\(\varphi_E:\Phi_{\Spec\bar k}^*E\xrightarrow{\sim}E\) give
\[
 \beta_Z:\Phi_Z^*\omega_Z
 \xrightarrow{\sim}Ra_Z^!\Phi_{\Spec\bar k}^*E
 \xrightarrow{Ra_Z^!(\varphi_E)}Ra_Z^!E=\omega_Z.
\]
This arrow is compatible with iteration, biduality, direct-image exchange,
and finite scalar extension.
\end{theorem}

\begin{theorem}[Strict Verdier duality]
Retain the preceding data and hypotheses.  For a strict Weil
complex \((K,\varphi_K)\), the strict Weil arrow on \(D_ZK\) is the typed
composite
\[
\begin{aligned}
 \varphi_{D_ZK}:\Phi_Z^*D_ZK
 &\xrightarrow{\sim}
 R\mathcal Hom_E(\Phi_Z^*K,\Phi_Z^*\omega_Z)\\
 &\xrightarrow{R\mathcal Hom_E(\varphi_K^{-1},\beta_Z)}
 R\mathcal Hom_E(K,\omega_Z)=D_ZK.
\end{aligned}
\]
It defines an exact anti-equivalence
\[
 D_Z:\DW(Z_0,E)^{\mathrm{op}}\xrightarrow{\sim}\DW(Z_0,E)
\]
with strict biduality.  The two direct-image exchange isomorphisms are
isomorphisms in the strict Weil categories.  Every definition and theorem
below involving \(D_Z\), lower weights, or general exceptional direct-image
duality is conditional on the finite-level exceptional-duality hypotheses.
The finite-coefficient middle-extension pairing for curves stated later has
its own independent hypotheses.
\end{theorem}

\section{Weights and their exact calculus}

\begin{definition}[Punctual purity and mixedness]
Fix \(w\in\mathbf R\).  A constructible Weil sheaf \(\mathcal F\) is
punctually \(\iota\)-pure of weight \(w\) if the finite
\(q_x\)-Frobenius module \((\mathcal F_{\bar x},F_x,q_x)\) is
\(\iota\)-pure of weight \(w\) for every closed point \(x\).  It is mixed
of \(\iota\)-weights at most \(w\) if, after a finite extension of \(E\) in
\(\Qlb\), it has a finite filtration by Weil subsheaves whose nonzero
graded pieces are punctually \(\iota\)-pure of real weights \(u\leq w\).
Mixedness of weights at least \(w\) is defined using pure graded weights
\(u\geq w\).  Punctual purity at all conjugates is defined by imposing the
all-conjugates predicate on every closed-point Frobenius module.
\end{definition}

\begin{definition}[Upper and lower weights of complexes]
For \(K\in\DW(X_0,E)\) and \(w\in\mathbf R\), define
\[
 \WtLe(K,w)
 \quad\Longleftrightarrow\quad
 \mathcal H^m(K)\text{ is mixed of weights at most }w+m
 \text{ for every }m.
\]
If \(X_0\) is separated over \(k\), assume the finite-level
exceptional-duality hypotheses and use the strict Verdier duality above to
define
\[
                 \WtGe(K,w)\quad\Longleftrightarrow\quad
                 \WtLe(D_XK,-w).
\]
The complex is pure of weight \(w\) when both conditions hold.
\end{definition}

\begin{definition}[Integral mixedness and integral complex weights]
Fix \(w\in\mathbf R\).  A constructible Weil sheaf \(\mathcal F\) is
\emph{integrally mixed of
weights at most \(w\)} if, after a finite extension of \(E\), it has a
finite filtration by Weil subsheaves whose nonzero graded pieces are
punctually \(\iota\)-pure of weights
\[
                              u\in\mathbf Z,\qquad u\leq w.
\]
For \(K\in\DW(X_0,E)\), define
\[
 \IntWtLe(K,w)
 \quad\Longleftrightarrow\quad
 \mathcal H^m(K)\text{ is integrally mixed of weights at most }w+m
 \text{ for every }m.
\]
When \(X_0\) is separated over \(k\), define
\[
                 \IntWtGe(K,w)
       \quad\Longleftrightarrow\quad
                 \IntWtLe(D_XK,-w).
\]
These conditions imply \(\WtLe(K,w)\) and \(\WtGe(K,w)\), respectively.
The superscript \(\mathbf Z\) constrains the weights of the pure graded
pieces; the numerical bound \(w\) itself may be real.
\end{definition}

\begin{proposition}[Exact calculus of upper weights]
Fix \(w\in\mathbf R\).
Let \(\mathsf M_{\le w}(X_0,E)\) be the full subcategory of constructible
Weil sheaves mixed of weights at most \(w\).
\begin{enumerate}[label=\textup{(\roman*)}]
\item \(\mathsf M_{\le w}\) is closed under Weil subobjects, quotients,
extensions, finite direct sums, and direct summands.
\item If \(K'\to K\to K''\to K'[1]\) is distinguished and
\(\WtLe(K',w)\) and \(\WtLe(K'',w)\), then \(\WtLe(K,w)\).
\item \(\WtLe(K[r],w+r)\) is equivalent to \(\WtLe(K,w)\).
\item The standard truncation triangle
\[
 \tau_{\le m-1}K\longrightarrow\tau_{\le m}K
 \longrightarrow\mathcal H^m(K)[-m]\longrightarrow
\]
is Frobenius equivariant.  Consequently, \(\WtLe(K,w)\) is equivalent to
the collection of bounds
\(\mathcal H^m(K)\in\mathsf M_{\le w+m}\).
\end{enumerate}
All assertions remain valid after finite scalar extension.
They also remain valid with \(\WtLe\) replaced by \(\IntWtLe\): the pure
weights occurring in subobjects, quotients, extensions, finite direct
sums, direct summands, shifts, and truncation towers remain integral.
\end{proposition}

\begin{proposition}[Local detection on a finite stratification]
Fix \(w\in\mathbf R\).  Let
\(X_0=\bigsqcup_{a=1}^rS_a\) be a finite stratification by locally
closed subschemes and let \(s_a:S_a\hookrightarrow X_0\).  A constructible
Weil sheaf \(\mathcal F\), already known to be mixed, has weights at most
\(w\) if and only if every \(s_a^*\mathcal F\) does.  Equivalently, after
refining so that all graded pieces are lisse, the bound is detected by the
closed-point Frobenius eigenvalues on all strata.  A complex \(K\) satisfies
\(\WtLe(K,w)\) if and only if these tests hold for every
\(s_a^*\mathcal H^m(K)\) with bound \(w+m\).
\end{proposition}

\begin{proposition}[Finite extension of the ground field]
For \(k_r=\mathbf F_{q^r}\), pullback from \(k\) to \(k_r\) replaces each
local Frobenius by a positive power and replaces the residue-field cardinality
by the same power.  Punctual purity and both weight inequalities are
therefore invariant under finite extension of the ground field.  They may
be checked after such an extension.
\end{proposition}

\section{Stratification and derived devissage}

\begin{definition}[Ordered support stratification]
An ordered support stratification of a noetherian scheme \(X_0\) is a chain
of closed reduced subschemes
\[
 \varnothing=X_{-1}\subset X_0'\subset\cdots\subset X_r'=X_0
\]
such that \(S_a=X_a'\setminus X_{a-1}'\) is smooth and locally closed in
\(X_0\).  Write \(s_a:S_a\hookrightarrow X_0\) for the locally closed
immersion.  The ordering ensures that adjoining one stratum is an
open-closed decomposition of \(X_a'\).
\end{definition}

\begin{theorem}[Pure-lisse stratification of a mixed sheaf]
Fix \(w\in\mathbf R\).  Let \(\mathcal F\) be an integrally mixed constructible Weil sheaf of
weights at most \(w\).  After a finite extension \(M/E\), there are an ordered support
stratification and a finite Frobenius-stable filtration of
\(\mathcal F_M\) whose successive quotients have the form
\[
                         s_{a,!}\mathcal L_a,
\]
where \(\mathcal L_a\) is lisse and punctually \(\iota\)-pure of an integral weight
\(u_a\in\mathbf Z\), \(u_a\le w\), on the smooth stratum \(S_a\).  A further refinement may be
chosen simultaneously for finitely many sheaves, morphisms, kernels, and
cokernels.
\dep{Noetherian constructible stratification and the exact calculus of
mixed Weil sheaves.}
\end{theorem}

\begin{theorem}[Sheaf devissage]
Fix \(f_0:X_0\to Y_0\) and \(w\in\mathbf R\).  Suppose that for every smooth
locally closed \(s:S_0\hookrightarrow X_0\) and every lisse punctually \(\iota\)-pure
Weil sheaf \(\mathcal L\) of integral weight
\(u\in\mathbf Z\), \(u\le w\),
\[
                     R^m(f_0s)_!\mathcal L
       \text{ is integrally mixed of weights at most }u+m
\]
for all \(m\).  Then the same conclusion with bound \(w+m\) holds for every
integrally mixed constructible Weil sheaf on \(X_0\) of weights at most
\(w\); moreover, every output sheaf is integrally mixed.
\dep{Pure-lisse stratification, exactness of \(s_!\) for a locally closed
immersion, composition of exceptional direct image, and extension closure
of upper weights.}
\end{theorem}

\begin{theorem}[Bounded-complex devissage]
Let \(f_0:X_0\to Y_0\) be separated and of finite type.  Assume that for
every \(a\in\mathbf R\), every integrally mixed constructible Weil sheaf
\(\mathcal F\) of weights at most \(a\), and every \(r\in\mathbf Z\),
\[
 R^rf_!\mathcal F\text{ is integrally mixed of weights at most }a+r.
\]
Then, for every \(w\in\mathbf R\) and every \(K\in\DW(X_0,E)\),
\(\IntWtLe(K,w)\) implies
\[
                         \IntWtLe(Rf_{!,\mathrm W}K,w).
\]
For integers \(m\leq n\) containing the cohomological support of \(K\), the
finite standard truncation tower has factors
\(\mathcal H^b(K)[-b]\), \(m\leq b\leq n\), and its functorial
hypercohomology spectral sequence is
\[
 E_2^{a,b}=R^af_!\mathcal H^b(K)
 \quad\Longrightarrow\quad
 \mathcal H^{a+b}(Rf_{!,\mathrm W}K).
\]
This spectral sequence is bounded, Frobenius equivariant, and compatible
with finite scalar extension.  Each abutment has a finite Frobenius-stable
filtration whose graded terms are subquotients of the corresponding
\(E_2^{a,b}\).
\dep{Bounded constructibility, the finite standard truncation tower, the
hypercohomology spectral sequence, and closure of integral upper weights
under subquotients and extensions.}
\end{theorem}

\begin{theorem}[Boundary-separated lisse devissage]
Let \(C_0\) be a smooth curve, let \(j:W_0\hookrightarrow C_0\) be a dense
open immersion, and fix a finite set \(B_0\subset W_0\) disjoint from
\(C_0\setminus W_0\).  Let \(\mathcal F\) be a lisse Weil sheaf on \(W_0\)
that is punctually \(\iota\)-pure of an integral weight
\(w\in\mathbf Z\) and unramified at every point of \(B_0\).  After one
finite extension of \(E\), there is a finite filtration by lisse Weil
subsheaves
\[
 0=\mathcal F_0\subset\mathcal F_1\subset\cdots\subset
 \mathcal F_N=\mathcal F
\]
whose quotients are irreducible lisse Weil sheaves, punctually
\(\iota\)-pure of weight \(w\), and unramified at every point of \(B_0\).
Each quotient is either geometrically constant or has no geometrically
constant simple subquotient.  The filtration and this dichotomy use only
the noetherian lisse Weil category, closed-point Frobenius modules, and
inertia invariants; no direct-image weight theorem is assumed.
\dep{Finite length of the lisse Weil category after finite scalar
extension, exactness of stalks, exactness of inertia invariants for the
trivial inertia action at \(B_0\), and fixed-\(\iota\) purity of
subquotients of a pure Frobenius module.}
\end{theorem}

\subsection{Localization with the correct bounds}

Let \(j:U_0\hookrightarrow X_0\) be open and let
\(i:Z_0\hookrightarrow X_0\) be its closed complement.

\begin{theorem}[Upper localization triangle]
For every \(w\in\mathbf R\) and \(K\in\DW(X_0,E)\), there is a
Frobenius-equivariant distinguished
triangle
\[
 j_!j^*K\longrightarrow K\longrightarrow i_*i^*K
                  \longrightarrow j_!j^*K[1].
\]
The functors \(j^*,i^*,j_!\), and \(i_*=i_!\) preserve upper weights without
changing the bound.  Consequently,
\[
 \WtLe(K,w)\quad\Longleftrightarrow\quad
 \WtLe(j^*K,w)\text{ and }\WtLe(i^*K,w).
\]
For \(s=j\) or \(i\), the comparison
\(Rf_!s_!\simeq R(fs)_!\) identifies the localization long exact sequence
with the one used in source devissage.
The same equivalence holds with \(\WtLe\) replaced by \(\IntWtLe\).
\end{theorem}

\begin{theorem}[Lower localization triangle]
Assume the finite-level exceptional-duality hypotheses and suppose that
\(X_0\) is separated over \(k\).  For every \(w\in\mathbf R\) and
\(K\in\DW(X_0,E)\), there is a Frobenius-equivariant
distinguished triangle
\[
 i_*Ri^!K\longrightarrow K\longrightarrow Rj_*j^*K
                  \longrightarrow i_*Ri^!K[1].
\]
The functors \(Ri^!,j^!,i_*\), and \(Rj_*\) preserve lower weights without
changing the bound.  These assertions are the Verdier duals of the upper
localization assertions.
\end{theorem}

\section{Finite maps, radicial maps, targets, and composition}

\begin{theorem}[Finite direct image]
Let \(u_0:T_0\to S_0\) be finite.  Then
\[
 Ru_{!,\mathrm W}=Ru_{*,\mathrm W}=u_*:
 \DW(T_0,E)\longrightarrow\DW(S_0,E),
\]
and \(u_*\) is exact on constructible Weil sheaves.  For a geometric point
\(\bar s\to S_0\),
\[
             (u_*\mathcal F)_{\bar s}
       \simeq\bigoplus_{\bar t\in T_{\bar s}}\mathcal F_{\bar t}.
\]
If \(s\) is closed, iterate its local geometric Frobenius through an orbit
of length \(d\).  The resulting action on one summand is the local
geometric Frobenius at the corresponding point of residue degree \(d\).
Hence \(u_*\) preserves punctual purity, upper weights, and lower weights.
It also preserves integral mixedness and the conditions \(\IntWtLe\) and
\(\IntWtGe\).
\end{theorem}

\begin{theorem}[Separated quasi-finite reduction]
Let \(u_0:T_0\to S_0\) be separated and quasi-finite.  There is a finite
ordered stratification \(T_0=\bigsqcup T_a\) and factorizations
\[
        T_a\xrightarrow{j_a}\overline T_a\xrightarrow{\bar u_a}S_0
\]
with \(j_a\) open and \(\bar u_a\) finite.  Therefore \(Ru_!\) preserves
ordinary and integral upper weights.  On sheaves it is concentrated in
degree zero.
\dep{Zariski's Main Theorem, finite direct image, and upper localization.}
\end{theorem}

\begin{theorem}[Radicial invariance]
Let \(v_0:S_0'\to S_0\) be a finite universal homeomorphism.  Pullback and
pushforward induce mutually quasi-inverse exact equivalences
\[
 \DW(S_0,E)\simeq\DW(S_0',E).
\]
They identify geometric stalks and the corresponding powers of geometric
Frobenius, commute with \(Rf_!\) after base change, and preserve punctual
purity and both weight inequalities.  In particular, a finite purely
inseparable extension of a generic ground field does not change a weight
statement.
\end{theorem}

\begin{theorem}[Target open-closed detection]
Fix \(w\in\mathbf R\) and \(K\in\DW(X_0,E)\).  Consider a Cartesian base
change of \(f_0:X_0\to Y_0\) along an open
immersion \(a:V_0\hookrightarrow Y_0\) and its closed complement
\(b:T_0\hookrightarrow Y_0\).  Put
\(X_{V,0}=X_0\times_{Y_0}V_0\) and
\(X_{T,0}=X_0\times_{Y_0}T_0\), and denote the projections to \(X_0\) by
\(a'\) and \(b'\), and the induced maps to \(V_0\) and \(T_0\) by
\(f_V\) and \(f_T\).  Compact-support base change gives
\[
 a^*Rf_!K\xrightarrow{\sim}Rf_{V,!}a'^*K,
 \qquad
 b^*Rf_!K\xrightarrow{\sim}Rf_{T,!}b'^*K.
\]
If both right-hand sides have upper complex weight \(w\), then so does
\(Rf_!K\).  The same assertion holds for a finite ordered stratification of
the target.
The statement remains valid with \(\WtLe\) replaced by \(\IntWtLe\).
\dep{Compact-support base change and local detection on a finite
stratification.}
\end{theorem}

\begin{theorem}[Leray weight transfer]
Let \(X_0\xrightarrow h Z_0\xrightarrow gY_0\) be separated and of finite
type, let \(w\in\mathbf R\), and let \(\mathcal F\) be a constructible
Weil sheaf.  The bounded
Frobenius-equivariant spectral sequence
\[
 E_2^{a,b}=R^ag_!R^bh_!\mathcal F
       \Longrightarrow R^{a+b}(gh)_!\mathcal F
\]
has a finite Frobenius-stable filtration on the abutment with graded pieces
\(E_\infty^{a,b}\).  If \(R^bh_!\mathcal F\) is integrally mixed of
weights at most \(w+b\) and \(R^ag_!(R^bh_!\mathcal F)\) is integrally
mixed of weights at most \(w+a+b\), then
\[
             R^n(gh)_!\mathcal F
       \text{ is integrally mixed of weights at most }w+n.
\]
\dep{The composition isomorphism for \(Rf_!\), boundedness, and the
subquotient and extension closure of upper weights.}
\end{theorem}

\section{From the affine line to the projective-curve theorem}

\subsection{The affine-line input}

Let \(A_0=\mathbf A^1_k\), let \(j:U_0\hookrightarrow A_0\) be dense, and
write \(A,U\) for the geometric base changes.  For a lisse Weil sheaf
\(\mathcal L\) on \(U_0\), the notation \(j_*\mathcal L\) in this section
means the underived direct-image sheaf.

Let \(A'_0\) be the dual affine line, let
\(p:A\times A'\to A\), \(p':A\times A'\to A'\) be the projections, and let
\(\mu:A\times A'\to\mathbf A^1_{\bar k}\) be \((x,x')\mapsto xx'\).
The Artin--Schreier sheaf \(\mathcal L_\psi\) is the rank-one lisse Weil
sheaf associated with the finite etale \((k,+)\)-torsor
\[
                              t^q-t=x
\]
and the character \(\psi:k\to\Qlb^\times\).

\begin{definition}[Fourier--Deligne transform and its local module]
For \(K\in\DW(A_0,E)\), define the normalized Fourier--Deligne transform
by
\[
 \operatorname{FT}_\psi(K)
   =Rp'_!\bigl(p^*K\otimes_E^L\mu^*\mathcal L_\psi\bigr)[1]
   \in\DW(A'_0,E).
\]
Its strict Frobenius arrow is the composite formed from the strict arrows on
\(K\) and \(\mathcal L_\psi\), the tensor arrow, the pullback exchanges,
and the compact-support exchange for \(p'\).  If the restriction of
\(\mathcal G=\operatorname{FT}_\psi(K)\) to a punctured strict henselian
neighborhood of \(0'\in A'_0\) is lisse and tamely ramified, its
\emph{local Fourier module at \(0'\)} is the finite \(q\)-Frobenius module
\((V',F,q)\), together with the nilpotent logarithm \(\widetilde N\) of
tame inertia.  The relation
\[
                  F\widetilde NF^{-1}=q^{-1}\widetilde N
\]
is equivalently a Frobenius-equivariant arrow \(N:V'\to V'(-1)\).
\end{definition}

\begin{theorem}[Fourier constructibility, base change, and inversion]
The Fourier--Deligne transform is an exact functor of ordinary triangulated
strict Weil categories and preserves bounded constructibility.  It commutes
with finite scalar extension and finite extension of the ground field.  If
\(a:A\to A\) is \(x\mapsto-x\), there is a Frobenius-equivariant inversion
isomorphism
\[
       \operatorname{FT}_{\psi^{-1}}\operatorname{FT}_\psi(K)
                  \xrightarrow{\sim}a^*K(-1).
\]
The local module and its inertia invariants commute with these coefficient
and ground-field changes.
\dep{Bounded constructibility and compact-support base change for \(Rp'_!\),
the strict tensor and pullback realizations, the Artin--Schreier convolution
calculation, and the Tate normalization.}
\end{theorem}

\begin{definition}[Fixed-\(\iota\) type-3 affine-line sheaf]
Fix \(w\in\mathbf R\).  A fixed-\(\iota\) type-3 sheaf of weight \(w\) is
a lisse Weil sheaf \(\mathcal L\) on
\(U_0\) which is irreducible in
\(\mathsf{Cons}_{\mathrm W}(U_0,E)\), punctually \(\iota\)-pure of weight
\(w\), and unramified at \(\infty\in\mathbf P^1_k\).  In addition, after
extension to \(\Qlb\), its geometric local system has no simple subquotient
which, after descent to a finite extension \(k'/k\), is the restriction of
an Artin--Schreier sheaf
\[
             \mathcal L_{\psi\circ\operatorname{Tr}_{k'/k}}(ax),
             \qquad a\in k',
\]
for a fixed nontrivial additive character \(\psi:k\to\Qlb^\times\).
Unramified at infinity means that geometric inertia at infinity acts
trivially on the geometric generic stalk.  Under this unramifiedness
hypothesis, the last condition is equivalent to the absence of a
geometrically constant simple subquotient; geometric irreducibility is not
required.
\end{definition}

\begin{definition}[All-conjugates type-3 affine-line sheaf]
An all-conjugates type-3 sheaf of weight \(w\in\mathbf R\) is a
fixed-\(\iota\) type-3 sheaf whose closed-point Frobenius modules are pure
of weight \(w\) at all conjugates.  This is a stronger coefficient
condition.  Below, the unqualified expression ``type-3 sheaf'' means a
fixed-\(\iota\) type-3 sheaf; an all-conjugates hypothesis is always stated
explicitly.
\end{definition}

\begin{theorem}[Affine-line numerical upper-weight estimate]
For every fixed-\(\iota\) type-3 sheaf \(\mathcal L\) of integral weight
\(w\in\mathbf Z\),
\[
 H_c^r(A,j_*\mathcal L)=0\quad(r\ne1),
 \qquad
 |\iota(\alpha)|\le q^{(w+1)/2}
\]
for every eigenvalue \(\alpha\) of geometric Frobenius on
\(H_c^1(A,j_*\mathcal L)\).
The Frobenius representation and its bound commute with finite extension of
the coefficient field and of the finite ground field.
\dep{Fourier--Deligne transformation and inversion, stationary phase,
local Fourier transformation, type-3 monodromy, and the fixed-\(\iota\)
local Fourier weight comparison.}
\end{theorem}

\begin{theorem}[Integral local-Fourier filtration]
Let \(\mathcal L\) be a fixed-\(\iota\) type-3 sheaf of integral weight
\(w\in\mathbf Z\), and let \(\mathcal G\) be its normalized
Fourier--Deligne transform.  Let \(V'\) be its tame local Frobenius module
at the distinguished dual point \(0'\).  After choosing a basis of the
one-dimensional Tate line \(E(1)\), let
\(\widetilde N:V'\to V'\) be the resulting nilpotent logarithm of tame
inertia, satisfying \(F\widetilde NF^{-1}=q^{-1}\widetilde N\); equivalently
it underlies a Frobenius-equivariant arrow \(N:V'\to V'(-1)\).  Put
\[
                  B'=\operatorname{im}\widetilde N,
        \qquad    K'=\ker\widetilde N.
\]
Changing the Tate basis rescales \(\widetilde N\) by a nonzero scalar and
does not change \(B'\), \(K'\), or the filtration below.
Then \(\widetilde N^2=0\), equivalently \(N(-1)\circ N=0\), and the
monodromy filtration centered at zero is
\[
 0\subset M_{-1}=B'\subset M_0=K'\subset M_1=V'.
\]
 After one finite extension \(E'/E\),
\begin{align*}
 \operatorname{gr}_{-1}^M=B'&\text{ is \(\iota\)-pure of weight }w,\\
 \operatorname{gr}_{0}^M=K'/B'&\text{ is \(\iota\)-pure of weight }w+1,\\
 \operatorname{gr}_{1}^M=V'/K'&\text{ is \(\iota\)-pure of weight }w+2.
\end{align*}
and \(N:\operatorname{gr}_1^M\xrightarrow{\sim}
\operatorname{gr}_{-1}^M(-1)\).  Stationary phase gives a
Frobenius-equivariant isomorphism
\[
 H_c^1(A,j_*\mathcal L)\otimes_EE'\xrightarrow{\sim}K'.
\]
Consequently \(0\subset B'\subset K'\) is a Frobenius-stable filtration
of the compactly supported cohomology with pure integral graded weights
\(w\) and \(w+1\).
\dep{Integral-weight Fourier purity for \(\operatorname{gr}_1^M\), the
length-two monodromy theorem at \(0'\),
\(F\widetilde NF^{-1}=q^{-1}\widetilde N\), the dual
pairing on \(\operatorname{gr}_0^M\), the boundary-invariant upper
estimates for the local module and its dual, and the stationary-phase
isomorphism \(H_c^1\simeq K'\).}
\end{theorem}

\begin{theorem}[Affine-line upper-weight estimate, integral-weight form]
For every fixed-\(\iota\) type-3 sheaf \(\mathcal L\) of integral weight
\(w\in\mathbf Z\),
\[
 H_c^r(A,j_*\mathcal L)=0\quad(r\ne1),
 \qquad
 H_c^1(A,j_*\mathcal L)
       \text{ is integrally mixed of weights at most }w+1.
\]
The Frobenius representation and its bound commute with finite extension of
the coefficient field and of the finite ground field.
\dep{The affine-line numerical upper-weight estimate, the integral
local-Fourier filtration, and the definition of integral mixedness for a
finite Frobenius module.}
\end{theorem}

\begin{definition}[Additional real-weight affine-line hypotheses]
Let \(\mathcal L\) be a fixed-\(\iota\) type-3 sheaf of real weight \(w\).
A \emph{boundary-real conjugate companion} is a fixed-\(\iota\) type-3 sheaf
\(\mathcal L^c\), after one finite coefficient extension, such that the
closed-point Frobenius spectra of \(\mathcal L^c\) are carried by \(\iota\)
to the complex conjugates of those of \(\mathcal L\), and the same equality
holds for the characteristic polynomials on every finite-boundary
inertia-invariant space.  The \emph{Fourier-output purity hypothesis}
asserts that the lisse restriction of the normalized Fourier--Deligne
transform of \(\mathcal L\) is punctually \(\iota\)-pure of its predicted
weight \(w+1\).
\end{definition}

\begin{theorem}[Conditional real-weight affine-line estimate]
Let \(\mathcal L\) be a fixed-\(\iota\) type-3 sheaf of real weight \(w\).
If either a
boundary-real conjugate companion exists or the Fourier-output purity
hypothesis holds, then
\[
 H_c^r(A,j_*\mathcal L)=0\quad(r\ne1),
 \qquad
 H_c^1(A,j_*\mathcal L)
       \text{ has \(\iota\)-weights at most }w+1.
\]
\dep{Fourier--Deligne inversion and stationary phase; the boundary-polynomial
comparison in the companion case; and the type-3 Fourier weight comparison
in the Fourier-output case.}
\end{theorem}

\subsection{The curve link}

Let \(C_0\) be a smooth projective geometrically connected curve,
\(j:U_0\hookrightarrow C_0\) a dense open immersion, and
\(\mathcal L\) an irreducible lisse Weil sheaf on \(U_0\).

\begin{proposition}[Good-infinity projection]
After a finite extension \(k'/k\), there is a finite generically etale
morphism
\[
                         \pi:C_{0,k'}\longrightarrow\mathbf P^1_{k'}
\]
such that \(\pi^{-1}(\infty)\) is a nonempty reduced divisor contained in
\(U_{0,k'}\), and \(\pi\) is etale over an open neighborhood of infinity.
The point at infinity is \(k'\)-rational.  Consequently, a lisse sheaf on
\(U_{0,k'}\) is lisse at every point over infinity, and its finite direct
image is lisse on an open neighborhood of infinity.
\dep{Riemann--Roch on \(C_0\), a reduced pole divisor supported in \(U_0\),
and the etale-locus theorem for a finite morphism of smooth curves.}
\end{proposition}

\begin{theorem}[Punctual correction upper bound]
Retain the good-infinity projection and suppose that \(\mathcal L\) is
punctually \(\iota\)-pure of an integral weight \(w\).  For
\(\mathcal G=\pi_*j_*\mathcal L\) and the maximal lisse open
\(j':W_0\hookrightarrow\mathbf P^1_{k'}\), the kernel and cokernel of
\[
                 \mathcal G\longrightarrow j'_*j'^*\mathcal G
\]
are punctual Weil sheaves.  After a finite coefficient extension, each of
their stalks has a finite Frobenius-stable filtration with pure integral
graded weights at most \(w\).  The assertion is stable under subquotients,
finite direct sums, and further finite extension of the coefficient or
ground field.
\dep{The integral local-Fourier filtration and stationary-phase comparison
for the local inertia modules, together with exact-sequence stability of
integral upper weights.}
\end{theorem}

\begin{proposition}[Finite-pushforward comparison and constituent
devissage]
Suppose that \(\mathcal L\) is punctually \(\iota\)-pure of an integral
weight \(w\in\mathbf Z\), and choose a good-infinity projection.  Finite
direct image is exact and there
is a Frobenius-equivariant identification
\[
 R\Gamma(C,j_*\mathcal L)
 \simeq R\Gamma(\mathbf P^1,\mathcal G),
 \qquad \mathcal G=\pi_*j_*\mathcal L.
\]
Let \(j':W_0\hookrightarrow\mathbf P^1_{k'}\) be the maximal open on which
\(\mathcal G\) is lisse.  This open contains infinity.  The canonical map
\(\mathcal G\to j'_*j'^*\mathcal G\) has punctual kernel and cokernel.
After finite scalar extension, \(\mathcal G\) has a Frobenius-stable finite
filtration whose
graded pieces are either punctual sheaves at points of
\(\mathbf A^1_{k'}\), or middle extensions \(j'_*\mathcal M\) of
irreducible lisse Weil sheaves on \(W_0\).  Every such \(\mathcal M\) is
punctually \(\iota\)-pure of weight \(w\), and every punctual graded factor
is integrally mixed of weights at most \(w\).  The restriction of \(\mathcal M\) to
\(W_0\cap\mathbf A^1\) is unramified at infinity.
\dep{Good-infinity projection, exact finite direct image, the power-change
invariance of finite Frobenius-module weights, the punctual correction
upper bound, boundary-separated lisse
devissage with the point at infinity distinguished, constructible
Jordan--Hoelder devissage, and the kernel--cokernel decomposition of the
map from a constructible sheaf to the middle extension of its maximal
lisse restriction.}
\end{proposition}

\begin{proposition}[Affine-projective localization sequence]
Let \(a:\mathbf A^1\hookrightarrow\mathbf P^1\) have complement infinity,
and let \(\mathcal G\) be a constructible Weil sheaf of the form
\(j'_*\mathcal M\).  The triangle
\[
 a_!a^*\mathcal G\longrightarrow\mathcal G
 \longrightarrow i_*\mathcal G_\infty\longrightarrow
\]
gives a Frobenius-equivariant long exact sequence.  If \(\mathcal M\) is
irreducible and nonconstant, the relevant part is
\[
 0\longrightarrow\mathcal G_\infty
 \longrightarrow H_c^1(\mathbf A^1,a^*\mathcal G)
 \longrightarrow H^1(\mathbf P^1,\mathcal G)\longrightarrow0.
\]
For a geometrically constant factor \(M\),
\[
 H_c^r(\mathbf A^1,M)=0\ (r\ne2),\qquad
 H_c^2(\mathbf A^1,M)\simeq M(-1).
\]
\end{proposition}

\begin{proposition}[Punctual and geometrically constant constituents]
Retain the finite-pushforward situation and its integral weight
\(w\in\mathbf Z\).  Let \(\mathcal Q\) be a punctual graded factor of
\(\pi_*j_*\mathcal L\).  Then
\[
 H^r(\mathbf P^1,\mathcal Q)=0\quad(r\ne0),\qquad
 H^0(\mathbf P^1,\mathcal Q)=\bigoplus_x\mathcal Q_{\bar x},
\]
and the weight bound is the punctual bound on these stalks.  If
\(j'_*\mathcal M\) is a geometrically constant graded factor of weight
\(w\), it extends to a geometrically constant sheaf on \(\mathbf P^1\), and
\[
 H^0(\mathbf P^1,j'_*\mathcal M)=M,\qquad
 H^1(\mathbf P^1,j'_*\mathcal M)=0,\qquad
 H^2(\mathbf P^1,j'_*\mathcal M)=M(-1).
\]
For every such punctual or geometrically constant factor \(\mathcal Q\),
\[
 H^r(\mathbf P^1,\mathcal Q)
 \text{ is integrally mixed of weights at most }w+r.
\]
\dep{Cohomology of a punctual sheaf, the constant-sheaf calculation on
\(\mathbf P^1\), and the Tate-twist convention.}
\end{proposition}

\begin{proposition}[Nonconstant constituents are type-3]
Retain the finite-pushforward situation and its integral weight
\(w\in\mathbf Z\).  Let \(j'_*\mathcal M\) be a nonconstant irreducible
graded factor from the
finite-pushforward devissage, and restrict \(\mathcal M\) to
\(W_0\cap\mathbf A^1\).  Then \(\mathcal M\) is irreducible, punctually \(\iota\)-pure
of weight \(w\), geometrically nonconstant, and unramified at infinity.
It is therefore a fixed-\(\iota\) type-3 sheaf.  The affine-line estimate and the
affine-projective localization sequence give
\[
 H^r(\mathbf P^1,j'_*\mathcal M)
 \text{ integrally mixed of weights at most }w+r
\]
for every \(r\).  The same bound holds for the finite extension assembled
from all punctual, constant, and nonconstant graded factors.
\dep{Good-infinity projection, finite-pushforward constituent devissage,
the definition of type-3, the affine-line estimate, affine-projective
localization, and extension closure.}
\end{proposition}

\begin{theorem}[Finite-coefficient middle-extension Poincare pairing]
Let \(E/\Ql\) be finite, let \(C_0\) be a smooth projective geometrically
connected curve over a finite field \(k'\), let
\(j:U_0\hookrightarrow C_0\) be dense, and let \(\mathcal L\) be a finite
rank lisse \(E\)-Weil sheaf on \(U_0\).  For every \(r\in\mathbf Z\) there
is a canonical perfect Frobenius-equivariant pairing of finite
\(\#k'\)-Frobenius modules
\[
 H^r(C,j_*\mathcal L)\times
 H^{2-r}(C,j_*\mathcal L^\vee)
 \longrightarrow E(-1).
\]
It is functorial in \(\mathcal L\), compatible with finite extension of
\(E\) and \(k'\), and identifies the first factor with the dual of the
second factor twisted by \((-1)\).  Its hypotheses do not include the
finite-level exceptional-duality hypotheses for general finite-type
schemes.
\dep{Finite-torsion Poincare duality on a smooth curve, duality for the
middle extension \(j_*\mathcal L[1]\), a stable \(\mathcal O_E\)-lattice,
uniform finite cohomological amplitude, derived inverse limit, and
rationalization.}
\end{theorem}

\begin{proposition}[Dual completion of the reverse reduction]
Let \(C_0\) be a smooth projective geometrically connected curve, let
\(j:U_0\hookrightarrow C_0\) be dense, let \(\mathcal L\) be irreducible
and punctually \(\iota\)-pure of an integral weight \(w\in\mathbf Z\), and
let \(r\in\mathbf Z\).  The
finite-coefficient middle-extension
pairing is
\[
 H^r(C,j_*\mathcal L)\times
 H^{2-r}(C,j_*\mathcal L^\vee)
 \longrightarrow E(-1).
\]
An upper bound \(w+r\) for \(\mathcal L\), together with the corresponding
upper bound \(-w+2-r\) for \(\mathcal L^\vee\), forces purity of
\(H^r(C,j_*\mathcal L)\) of weight \(w+r\).
\dep{The nonconstant upper estimate applied separately to
\(\mathcal L\) and \(\mathcal L^\vee\), the punctual and constant
calculations, and the finite-coefficient middle-extension Poincare pairing
with its \((-1)\)-twist.}
\end{proposition}

\begin{theorem}[Affine-line-to-curve link]
If \(\mathcal L\) is irreducible and punctually \(\iota\)-pure of an
integral weight \(w\in\mathbf Z\), then
\[
                   H^r(C,j_*\mathcal L)
       \text{ is \(\iota\)-pure of the integral weight }w+r
\]
for every \(r\).  If \(\mathcal L\) is nonconstant and has no inertia
invariants at the missing points, then
\[
 H_c^0(U,\mathcal L)=H_c^2(U,\mathcal L)=0,
 \qquad H_c^1(U,\mathcal L)\text{ is \(\iota\)-pure of weight }w+1.
\]
\dep{Good-infinity projection, finite-pushforward comparison, punctual and
constant constituents, the type-3 nonconstant estimate, Frobenius-stable
coefficient devissage, dual completion, and the integral-weight affine-line
estimate after every finite coefficient and ground-field extension.}
\end{theorem}

\begin{theorem}[Absolute tame-boundary upper bound]
Let \(C_0\) be a smooth projective curve, let
\(j:U_0\hookrightarrow C_0\) be dense with reduced closed complement
\(i:D_0\hookrightarrow C_0\), and let \(\mathcal L\) be lisse and
punctually \(\iota\)-pure of an integral weight \(w\).  Then
\[
                         B=i^*j_*\mathcal L
\]
is a punctual Weil sheaf integrally mixed of weights at most \(w\).  The
short exact sequence
\[
 0\longrightarrow j_!\mathcal L\longrightarrow j_*\mathcal L
   \longrightarrow i_*B\longrightarrow0
\]
is Frobenius equivariant.  Together with the projective middle-extension
bound, it gives the compact-support upper bound for \(\mathcal L\) on
\(U_0\).
\dep{Quasi-unipotence of local monodromy, a finite tame cover and its trace
retraction, the punctual correction upper bound, local Fourier stationary
phase, and the open--closed localization sequence.}
\end{theorem}

\begin{corollary}[Integrally mixed constructible curve estimate]
Let \(a:C_0\to\Spec k\) be a separated curve and let
\(w\in\mathbf R\).  If \(K\in\DW(C_0,E)\) satisfies
\(\IntWtLe(K,w)\), then
\[
                         \IntWtLe(Ra_!K,w).
\]
\dep{Finite normalization and radicial invariance, pure-lisse
stratification, localization at finitely many points, the curve theorem,
the absolute tame-boundary upper bound, and complex devissage.}
\end{corollary}

\section{Controlled relative curves}

\subsection{Geometric data, stated without a weight conclusion}

\begin{definition}[Relative-curve compactification datum]
Let \(S_0\) be connected and normal.  A relative-curve compactification
datum over \(S_0\) consists of
\[
 X_0\xrightarrow{j}\overline X_0\xrightarrow{\bar f}S_0,
 \qquad D_0\xrightarrow{i}\overline X_0,
 \qquad g=\bar f i:D_0\to S_0,
\]
such that \(\bar f\) is smooth and projective of relative dimension one,
\(j\) is the open complement of \(i\), and \(g\) is finite etale.  The
original relative curve is \(f=\bar f j\).  This datum contains no assertion
about the weights of a direct image.
\end{definition}

\begin{definition}[Tame coefficient along the boundary]
Let \(\mathcal L\) be lisse on \(X_0\).  It is tame along \(D_0\) when every
wild inertia group at a geometric point of \(D_0\) acts trivially on its
geometric generic stalk.  Define the boundary invariant sheaf
\[
                         B=i^*j_*\mathcal L.
\]
Its geometric stalk is the invariant subspace of the corresponding local
inertia representation.  The symbol \(j_*\mathcal L\) is underived.
\end{definition}

\begin{theorem}[Generic controlled compactification]
Let \(h_0:Z_0\to S_0\) be a dominant separated finite-type morphism between
integral \(k\)-schemes whose generic fiber has dimension one and is
generically smooth, equivalently whose generic function-field extension is
separably generated.  Let
\(\mathcal L\) be lisse on a dense open of \(Z_0\).  There exist
\begin{enumerate}[label=\textup{(\roman*)}]
\item a dense open \(V_0\subseteq S_0\) and a dense open
      \(U_0\subseteq Z_0\times_{S_0}V_0\), with structure map
      \(h:U_0\to V_0\);
\item a finite radicial surjection \(v:V_0'\to V_0\), the Cartesian pullback
      \(U_{0,V'}=U_0\times_{V_0}V_0'\), its projection
      \(v_X:U_{0,V'}\to U_0\), and its structure map
      \(h':U_{0,V'}\to V_0'\);
\item a connected finite etale cover
      \(u:U_0'\to U_{0,V'}\) of constant degree \(d>0\);
\item a relative-curve compactification
      \(U_0'\xrightarrow j\overline U_0'\to V_0'\)
      with finite-etale boundary;
\item tameness of
      \(\mathcal L'=u^*v_X^*(\mathcal L|_{U_0})\) along that boundary.
\end{enumerate}
All schemes, sheaves, maps, adjunction morphisms, and compactification data
descend simultaneously from the generic point to the displayed dense open.
They remain compatible after a further shrinking of \(V_0\).
\dep{Generic smoothness, normalization of a curve, a finite cover killing
wild boundary monodromy, finite purely inseparable descent, relative
compactification of a smooth curve, and spreading out of finite data.}
\end{theorem}

\begin{theorem}[Trace retraction for the finite cover]
Use the notation of the generic controlled compactification.  On
\(U_{0,V'}\) there are Frobenius-equivariant maps
\[
 v_X^*\mathcal L\xrightarrow{\eta_u}
 u_*u^*v_X^*\mathcal L
 \xrightarrow{\operatorname{Tr}_u}v_X^*\mathcal L
\]
whose composite is \(d\,\mathrm{id}_{v_X^*\mathcal L}\).  Since \(E\) has
characteristic zero, \(d\) is invertible in \(E\), and
\(d^{-1}\operatorname{Tr}_u\) retracts \(\eta_u\).  Compact-support base
change gives
\[
 v^*R^rh_!\mathcal L
 \xrightarrow{\sim}R^rh'_!v_X^*\mathcal L.
\]
Under this isomorphism,
\[
 R^rh'_!v_X^*\mathcal L
 \quad\text{is a Frobenius-equivariant direct summand of}\quad
 R^r(h'u)_!u^*v_X^*\mathcal L.
\]
\dep{Finite-etale adjunction, the trace-degree formula, composition of
exceptional direct image, compact-support base change for the square
defined by \(v\), and radicial invariance.}
\end{theorem}

\subsection{Boundary and fiber theorems}

\begin{definition}[Local tame monodromy datum]
Fix a relative-curve compactification datum over \(S_0\), a lisse sheaf
\(\mathcal L\) on \(X_0\), and a closed point \(d\in D_0\) with image
\(s\in S_0\).  Put \(q_d=\#\kappa(d)\), choose a geometric point
\(\bar d\to d\), and let \(V\) be the local geometric generic stalk of
\(\mathcal L\), with local geometric Frobenius \(F_d\).  Let \(I\) be the
local tame inertia group and \(I_{\mathrm{fin}}\) its finite semisimple
quotient.  Choose a basis \(t_\ell\) of the one-dimensional Tate line and
write the unipotent part as
\[
                       T=\exp(t_\ell\widetilde N)
\]
for a nilpotent endomorphism \(\widetilde N\) of \(V\).  Its intrinsic form
is a morphism of finite \(q_d\)-Frobenius modules
\[
 N:(V,F_d,q_d)\longrightarrow(V,F_d,q_d)(-1).
\]
Equivalently,
\[
       F_d\widetilde NF_d^{-1}=q_d^{-1}\widetilde N.
\]
Changing the Tate basis rescales \(\widetilde N\) and does not change its kernel, image, or
monodromy filtration.
The monodromy filtration \(M_\bullet V\), centered at zero, is the unique
increasing filtration satisfying
\[
 N(M_aV)\subseteq M_{a-2}V(-1),\qquad
 N^a:\operatorname{Gr}_a^M V
       \xrightarrow{\sim}\operatorname{Gr}_{-a}^M V(-a)
       \quad(a\geq0).
\]
The relation with \(F_d\) and uniqueness of \(M_\bullet\) make every
\(M_aV\) Frobenius stable.
The inertia-invariant space is
\[
 V^I=(V^{I_{\mathrm{fin}}})\cap\ker N,
\]
where \(I_{\mathrm{fin}}\) denotes the finite semisimple part of tame
inertia.
\end{definition}

\begin{definition}[Monodromy-pure tame local datum]
A local tame monodromy datum
\((V,F_d,q_d;N,M_\bullet)\) is
\emph{monodromy-pure of weight \(w\in\mathbf R\)} as follows.  After a
finite coefficient extension splitting \(I_{\mathrm{fin}}\), geometric
Frobenius permutes its character eigenspaces.  For every Frobenius orbit
\(\mathcal O\) of characters, put
\(V_{\mathcal O}=\bigoplus_{\chi\in\mathcal O}V_\chi\).  The condition is
that every Frobenius-stable module
\[
 (\operatorname{Gr}_a^M V_{\mathcal O},F_d,q_d)
\]
is \(\iota\)-pure of weight \(w+a\) for every integer \(a\).
This is an additional local assertion: it does not follow merely from the
relations defining \(N\) and \(M_\bullet\).
\end{definition}

\begin{theorem}[Independent local-Fourier monodromy input]
Let \(j:U_0\hookrightarrow C_0\) be a dense open in a smooth curve over a
finite field, let \(\mathcal L\) be lisse and punctually
\(\iota\)-pure of an integral weight \(w\in\mathbf Z\), and assume that
\(\mathcal L\)
is tame at \(c\in C_0\setminus U_0\).  For the henselian local trait at
\(c\), let \((V,F_c,q_c;N,M_\bullet)\) be the local tame monodromy datum.
Then
this datum is monodromy-pure of weight \(w\).  The statement is unchanged
after a finite extension of the coefficient field or residue field and may
be applied to every closed geometric fiber of a smooth relative curve.
\dep{Exactness and base change of local Fourier transform, local Fourier
inversion, the local Fourier stationary-phase sequence, integral-weight
Fourier purity, type-3 length-two monodromy, type-3 weight comparison, and
finite semisimple-inertia eigenspace devissage.}
\end{theorem}

\begin{theorem}[Invariant upper bound from monodromy purity]
Let \((V,F_d,q_d;N,M_\bullet)\) be a monodromy-pure tame local datum of weight
\(w\in\mathbf Z\).  Then \(\ker N\) is integrally mixed of weights at
most \(w\), and
\[
                 V^I=(V^{I_{\mathrm{fin}}})\cap\ker N
\]
is a Frobenius-stable subspace integrally mixed of weights at most \(w\).
The same bound holds on every Frobenius-orbit sum of
\(I_{\mathrm{fin}}\)-character eigenspaces.  In particular, the trivial
character summand used in \(V^I\) is Frobenius stable.
\dep{Primitive decomposition for the monodromy filtration, the purity of
\(\operatorname{Gr}_a^M V\) of weight \(w+a\), and subobject and extension
closure of upper weights.}
\end{theorem}

\begin{theorem}[Spreading and fiber compatibility of monodromy]
Let a relative-curve compactification datum over \(S_0\) and a lisse
coefficient \(\mathcal L\) tame along \(D_0\) be given.  After shrinking
\(S_0\), the finite tame-inertia action, the nilpotent
operator \(N\), the subobjects \(M_a\), the graded objects
\(\operatorname{Gr}_a^M\), and the invariant object \(\ker N\) spread to
constructible sheaves along the finite-etale divisor \(D_0\).  Their
formation commutes, as an ordinary sheaf-theoretic comparison, with every
geometric fiber of \(S_0\).  If \(d\in D_0\) is a closed boundary point
over a closed point \(s\), the comparison intertwines local geometric
Frobenius after the \([\kappa(d):k]\)-fold descent iterate and preserves the
displayed relation with \(q_d^{-1}\).
\dep{Spreading out of a finite tame representation, uniqueness of the
monodromy filtration, constructibility, and proper/finite base change along
the boundary.}
\end{theorem}

\begin{theorem}[Relative tame boundary theorem]
Let a relative-curve compactification datum be given and let
\(\mathcal L\) be lisse, punctually \(\iota\)-pure of an integral weight
\(w\in\mathbf Z\), and tame along the boundary.  Then
\(B=i^*j_*\mathcal L\) is constructible and integrally mixed of weights
at most \(w\).  At a geometric boundary point its stalk is the invariant
space \(V^I\) from the local tame monodromy datum.  Formation of \(B\), its
induced filtration, and every graded quotient commutes ordinarily with
arbitrary geometric fibers; over closed points these comparisons commute
with the iterated local geometric Frobenius.
\dep{The independent local-Fourier monodromy input, the invariant upper
bound from monodromy purity, spreading and fiber compatibility, and
finite-etale boundary descent.}
\end{theorem}

\begin{theorem}[Relative boundary localization]
For the same datum there is a short exact sequence of constructible Weil
sheaves on \(\overline X_0\),
\[
 0\longrightarrow j_!\mathcal L\longrightarrow j_*\mathcal L
   \longrightarrow i_*B\longrightarrow0,
\]
and hence a Frobenius-equivariant distinguished triangle after applying
\(R\bar f_*\).  Put
\[
 A^r=R^rf_!\mathcal L,\qquad
 P^r=R^r\bar f_*j_*\mathcal L,\qquad
 E_B=g_*B.
\]
Finite direct image makes \(R^ag_*B=0\) for \(a\ne0\), and relative-curve
cohomological amplitude gives \(A^r=P^r=0\) outside \(0\le r\le2\).  The
long exact sequence therefore contains
\[
 0\to A^0\to P^0\to E_B\to A^1\to P^1\to0,
 \qquad A^2\xrightarrow{\sim}P^2.
\]
The triangle and long exact sequence are morphisms in the Weil category.
Their ordinary base-change comparisons are isomorphisms at every geometric
point.  Over a closed point \(s\), these isomorphisms intertwine local
geometric Frobenius after the \([\kappa(s):k]\)-fold descent iterate.
\dep{The boundary short exact sequence, proper base change, finite direct
image, and cohomological dimension two for relative curves.}
\end{theorem}

\begin{theorem}[Projective-fiber comparison]
Under the hypotheses of the relative tame boundary theorem, after shrinking
\(S_0\), each \(P^r\) is lisse, and for every geometric
point \(\bar s\to S_0\) proper base change gives an ordinary isomorphism
\[
 (P^r)_{\bar s}\xrightarrow{\sim}
 H^r(\overline X_{\bar s},j_{\bar s,*}\mathcal L_{\bar s}).
\]
If \(s\) is closed, both sides carry local geometric Frobenius and the
isomorphism is equivariant after the
\([\kappa(s):k]\)-fold descent iterate.
The affine-line-to-curve theorem applied to the right-hand side implies
that \(P^r\) is punctually \(\iota\)-pure of the integral weight \(w+r\).
\dep{Constructibility, generic lissity on the target, compatibility of
underived \(j_*\) with the controlled fibers, proper base change, and the
projective-curve theorem.}
\end{theorem}

\begin{theorem}[Controlled relative-curve estimate]
Under the hypotheses of the relative tame boundary theorem,
\[
                     R^rf_!\mathcal L
       \text{ is integrally mixed of weights at most }w+r
       \qquad(0\le r\le2).
\]
\dep{Relative tame boundary, relative boundary localization,
projective-fiber comparison, finite direct image, and exact weight
calculus.}
\end{theorem}

\begin{corollary}[Descent from the controlled cover]
In the situation of the generic controlled compactification, assume that
\(\mathcal L|_{U_0}\) is punctually \(\iota\)-pure of an integral weight
\(w\in\mathbf Z\) and that its pullback to the controlled cover is tame
along the finite-etale boundary.  For every \(0\leq r\leq2\), the controlled
relative-curve estimate for
\(u^*v_X^*(\mathcal L|_{U_0})\) makes
\(R^rh'_!v_X^*(\mathcal L|_{U_0})\) an integrally mixed direct summand of
weights at most \(w+r\).
The base-change isomorphism
\[
 v^*R^rh_!(\mathcal L|_{U_0})
 \xrightarrow{\sim}
 R^rh'_!v_X^*(\mathcal L|_{U_0})
\]
and radicial invariance show that
\(R^rh_!(\mathcal L|_{U_0})\) is integrally mixed of weights at most
\(w+r\) over \(V_0\).
\dep{Trace retraction, direct-summand closure, radicial invariance, and
compact-support base change.}
\end{corollary}

\section{Noetherian and support-dimension induction}

\begin{definition}[Support complexity]
For a nonempty constructible support \(Z\) in a finite-type scheme, define
\[
 \mu(Z)=\bigl(\dim Z,c(Z)\bigr)\in\mathbf N\times\mathbf N,
\]
where \(c(Z)\) is the number of irreducible components having maximal
dimension.  Order these pairs lexicographically.  If a dense open meets
every maximal-dimensional component, its closed complement has strictly
smaller complexity.  The empty support is the terminal case.
\end{definition}

\begin{theorem}[Finite support-dimension devissage]
Let \(\mathcal P(Z,\mathcal F)\) be a property of a constructible Weil sheaf
supported on \(Z\).  Assume:
\begin{enumerate}[label=\textup{(\roman*)}]
\item for every finite coefficient extension \(M/E\),
\(\mathcal P(Z,\mathcal F)\) holds if and only if
\(\mathcal P(Z,\mathcal F\otimes_E M)\) holds; moreover \(\mathcal P\) is
stable under subquotients, extensions, and direct summands;
\item for an open \(j:U\hookrightarrow Z\) with closed complement \(i:T\),
the two properties for \(j^*\mathcal F\) and \(i^*\mathcal F\) imply the
property for \(\mathcal F\);
\item on a dense smooth open of every irreducible support, \(\mathcal P\)
holds for every pure lisse constituent of integral weight.
\end{enumerate}
Then \(\mathcal P(Z,\mathcal F)\) holds for every integrally mixed
constructible Weil sheaf \(\mathcal F\).
\dep{Noetherian induction, pure-lisse stratification, and upper
localization.}
\end{theorem}

\begin{theorem}[Generic projection dichotomy]
Let \(q_0:S_0\to Y_0\) be the restriction to a smooth irreducible locally
closed subscheme \(S_0\subset\mathbf A^1_{Y_0}\).  After replacing the
scheme-theoretic image of \(S_0\) by a dense open, exactly one of the
following cases occurs.
\begin{enumerate}[label=\textup{(\roman*)}]
\item The generic fiber has dimension zero.  Then a dense open of \(S_0\)
is separated and quasi-finite over the target.
\item The generic fiber has dimension one.  Then a dense open of \(S_0\)
is a smooth relative curve, and every pure lisse coefficient of integral
weight on it admits
the generic controlled compactification and trace retraction of the
preceding section after the stated finite radicial base change.
\end{enumerate}
The removed source and target complements are proper closed subsets and
have strictly smaller support complexity.
\dep{The fiber-dimension theorem, generic flatness and smoothness,
normalization, the quasi-finite locus theorem, and generic controlled
compactification.}
\end{theorem}

\begin{theorem}[Relative affine-line projection]
Let \(Y_0\) be any finite-type \(k\)-scheme and let
\(p:\mathbf A^1_{Y_0}\to Y_0\) be the projection.  For every
\(w\in\mathbf R\) and every \(K\in\DW(\mathbf A^1_{Y_0},E)\),
\[
                    \IntWtLe(K,w)\quad\Longrightarrow\quad
                    \IntWtLe(Rp_{!,\mathrm W}K,w).
\]
For an integrally mixed sheaf \(\mathcal F\) of weights at most \(w\),
this says
\[
                R^rp_!\mathcal F
        \text{ is integrally mixed of weights at most }w+r.
\]
\dep{Complex and pure-lisse sheaf devissage; target open-closed detection;
support-dimension induction; the generic projection dichotomy; the
quasi-finite theorem in dimension zero; the controlled relative-curve
estimate in dimension one; trace, finite, and radicial descent.  The only
one-dimensional archimedean estimate entering the curve branch is the
affine-line theorem of the preceding section.}
\end{theorem}

\begin{proposition}[Local-to-global stalk detection for the projection]
Let \(Y_0\) be a finite-type \(k\)-scheme, let
\(p:\mathbf A^1_{Y_0}\to Y_0\) be the projection, let
\(w\in\mathbf R\), and let \(\mathcal F\) be an integrally mixed
constructible Weil sheaf on \(\mathbf A^1_{Y_0}\) of weights at most \(w\).
After the finite target stratification produced in the preceding theorem,
every \(R^rp_!\mathcal F\) is lisse on each target stratum.  Compact-support
base change identifies its stalk at a closed point \(y\) with
\[
               H_c^r(\mathbf A^1_{\bar y},\mathcal F|_{\mathbf A^1_{\bar y}}).
\]
The finite weight filtrations constructed in the quasi-finite and curve
branches restrict to the corresponding stalk filtrations.  Hence the
closed-point bounds \(w+r\) on all strata imply the global mixed bound on
\(R^rp_!\mathcal F\), with integral pure graded weights.
\dep{Constructibility, compact-support base change, finite target
stratification, and local detection for an already constructed mixed
filtration.}
\end{proposition}

\section{From projections to an arbitrary separated morphism}

\begin{proposition}[Affine generator factorization]
Let \(V_0=\Spec B\) be affine and let
\(f_0:U_0=\Spec A\to V_0\) be of finite type.  A finite set of \(B\)-algebra
generators of \(A\) gives a closed immersion
\[
 U_0\xrightarrow{c}\mathbf A^N_{V_0}.
\]
The projection \(\mathbf A^N_{V_0}\to V_0\) factors as \(N\) successive
relative affine-line projections.  The construction commutes with target
base change and geometric Frobenius.  Since a closed immersion is finite,
\(c_!=c_*\) is exact and weight-neutral.
More generally, the same factorization exists over a nonaffine target when
the \(\mathcal O_{V_0}\)-algebra \(f_{0,*}\mathcal O_{U_0}\) is generated by a
specified finite family of global sections.
\end{proposition}

\begin{theorem}[Affine direct image over an affine target]
Let \(V_0=\Spec B\) be affine and let
\(f_0:U_0=\Spec A\to V_0\) be of finite type.  For every
\(w\in\mathbf R\) and \(K\in\DW(U_0,E)\),
\[
 \IntWtLe(K,w)\quad\Longrightarrow\quad
 \IntWtLe(Rf_{!,\mathrm W}K,w).
\]
\dep{Affine generator factorization, finite direct image, repeated relative
affine-line projection, and Leray weight transfer.}
\end{theorem}

\begin{theorem}[Affine direct-image theorem]
Let \(f_0:U_0\to V_0\) be affine and of finite type, with \(V_0\) any
finite-type \(k\)-scheme.  Fix \(w\in\mathbf R\) and
\(K\in\DW(U_0,E)\).  Choose a finite affine open cover
\(V_0=\bigcup_{a=1}^mV_a\), and form the ordered locally closed affine
partition
\[
 T_a=V_a\setminus\bigcup_{b<a}V_b,
\]
discarding empty terms.  Each base change
\(U_0\times_{V_0}T_a\to T_a\) satisfies the affine-target theorem.
Compact-support base change and target open--closed detection therefore give
\[
 \IntWtLe(K,w)\quad\Longrightarrow\quad
 \IntWtLe(Rf_{!,\mathrm W}K,w)
\]
\dep{The affine-target theorem on every \(T_a\), integral upper
localization, compact-support base change, and finite induction on the
ordered target partition.}
\end{theorem}

\begin{theorem}[Finite affine refinement]
For every separated finite-type \(f_0:X_0\to Y_0\), there are finite ordered
open-closed stratifications \(Y_0=\bigsqcup V_a\) and
\(X_0\times_{Y_0}V_a=\bigsqcup U_{ab}\) such that every \(V_a\) and
\(U_{ab}\) is affine and every restricted map \(U_{ab}\to V_a\) is affine
of finite type.  The corresponding localization triangles and target
base-change isomorphisms are Frobenius equivariant.
\end{theorem}

\begin{theorem}[Global sheaf form over a finite coefficient field]
Let \(f_0:X_0\to Y_0\) be separated and of finite type.  If
\(w\in\mathbf R\) and \(\mathcal F\) is a constructible Weil sheaf
integrally mixed of weights
at most \(w\),
then
\[
                 R^if_!\mathcal F
       \text{ is integrally mixed of weights at most }w+i
\]
for every integer \(i\).
\dep{Finite affine refinement, the affine direct-image theorem, source
localization, target open-closed detection, and sheaf devissage.}
\end{theorem}

\begin{theorem}[Deligne's upper-weight theorem, integral-weight form]
Let \(f_0:X_0\to Y_0\) be separated and of finite type, let
\(K\in\DW(X_0,E)\), and let \(w\in\mathbf R\).  Then
\[
 \boxed{\qquad
      \IntWtLe(K,w)\quad\Longrightarrow\quad
      \IntWtLe(Rf_{!,\mathrm W}K,w).
 \qquad}
\]
Equivalently, \(\mathcal H^i(Rf_{!,\mathrm W}K)\) is integrally mixed of
weights at most \(w+i\) for every \(i\).
\dep{The global sheaf form and bounded-complex devissage.}
\end{theorem}

\begin{corollary}[Ordinary upper bound for integral inputs]
Under the hypotheses of the preceding theorem,
\[
          \IntWtLe(K,w)\quad\Longrightarrow\quad
          \WtLe(Rf_{!,\mathrm W}K,w).
\]
This is the ordinary upper-weight conclusion obtained by forgetting that
the output filtration has integral pure graded weights.
\end{corollary}

\begin{remark}[No unconditional real-weight conclusion]
The implication obtained by replacing \(\IntWtLe\) with \(\WtLe\) would
require the real-weight affine-line estimate for every type-3 constituent
appearing in the curve reductions and their duals.  The local-Fourier
module supplies that estimate only with its stated companion or
Fourier-output-purity hypothesis.  Accordingly, no theorem in this note
asserts that implication without one of those additional hypotheses.
\end{remark}

\section{Finite-stage algebraically closed coefficients}

\begin{definition}[Filtered finite-stage category]
Let \(\mathcal I\) be the filtered category of finite extensions
\(E/\Ql\) contained in \(\Qlb\).  First form the category whose objects are
pairs \((E,K_E)\), with \(K_E\in\DW(X_0,E)\), and whose morphisms are
\[
 \operatorname{Hom}((E,K_E),(E',L_{E'}))
 =\varinjlim_{M\supset E,E'}
 \operatorname{Hom}_{\DW(X_0,M)}
 (K_E\otimes_EM,L_{E'}\otimes_{E'}M).
\]
Declare a triangle distinguished when it is represented by a distinguished
triangle over one common finite stage, and take the Karoubi envelope.  The
resulting ordinary idempotent-complete triangulated category is denoted
\[
                       \DW(X_0,\Qlb)_{\mathrm{fin}}.
\]
An object may therefore be represented by a triple \((E,K_E,e_E)\), where
\(e_E^2=e_E\).
\end{definition}

\begin{definition}[Weights at an algebraically closed stage]
Fix \(w\in\mathbf R\).  For a representative \((E,K_E,e_E)\), define
\(\WtLe((E,K_E,e_E),w)\) to mean that the corresponding direct summand has
upper weights at most \(w\) after one, equivalently every sufficiently large, finite
coefficient extension.  Lower weights are defined by the stagewise Verdier
dual.  Scalar extension and direct-summand closure make these definitions
independent of the representative.  Define \(\IntWtLe\) and
\(\IntWtGe\) in the same way, requiring the representative to have a
finite filtration with integral pure graded weights at a finite stage.
\end{definition}

\begin{theorem}[Common finite-stage descent]
Fix a separated finite-type morphism \(f_0:X_0\to Y_0\).
Every finite collection of objects, idempotents, morphisms, distinguished
triangles, truncations, compactifications, Frobenius structures, and
six-operation comparison maps in
\(\DW(-,\Qlb)_{\mathrm{fin}}\) is represented over one finite
\(E/\Ql\).  Let \((E,K_E,e_E)\) be a representative of an object of
\(\DW(X_0,\Qlb)_{\mathrm{fin}}\), with \(e_E^2=e_E\).  For every finite
extension \(M/E\), there are coherent
isomorphisms
\[
 (Rf_{!,\mathrm W}K_E)\otimes_EM
 \xrightarrow{\sim}Rf_{!,\mathrm W}(K_E\otimes_EM),
\]
and similarly for \(Rf_*,D_X\), pullback, tensor product, and internal Hom.
They carry \(Rf_!(e_E)\) to the idempotent obtained after scalar extension.
\end{theorem}

\begin{theorem}[Algebraically closed coefficient form, integral weights]
Let \(f_0:X_0\to Y_0\) be separated and of finite type.
The stagewise functors descend to an exact functor
\[
 Rf_{!,\mathrm W}:\DW(X_0,\Qlb)_{\mathrm{fin}}
       \longrightarrow\DW(Y_0,\Qlb)_{\mathrm{fin}}.
\]
For every \(K\) in this category and every real \(w\),
\[
             \IntWtLe(K,w)\quad\Longrightarrow\quad
             \IntWtLe(Rf_{!,\mathrm W}K,w).
\]
\dep{Common finite-stage descent, the finite-coefficient theorem, and
closure of upper weights under direct summands.}
\end{theorem}

\section{Dual lower bounds and proper purity}

\begin{theorem}[Verdier exchange]
Assume the finite-level exceptional-duality hypotheses.  Let
\(f_0:X_0\to Y_0\) be separated and of finite type, and assume that
\(X_0\) and \(Y_0\) are separated over \(k\).  Verdier duality is an
exact anti-equivalence on the strict Weil categories and there are
Frobenius-equivariant natural isomorphisms
\[
 D_YRf_{*,\mathrm W}\xrightarrow{\sim}Rf_{!,\mathrm W}D_X,
 \qquad
 D_YRf_{!,\mathrm W}\xrightarrow{\sim}Rf_{*,\mathrm W}D_X.
\]
They commute with finite scalar extension and with the finite-stage
\(\Qlb\)-category.
\end{theorem}

\begin{theorem}[Lower-weight theorem for ordinary direct image,
integral-weight form]
Assume the finite-level exceptional-duality hypotheses.  Let
\(f_0:X_0\to Y_0\) be separated and of finite type with \(X_0\) and
\(Y_0\) separated over \(k\).  For every \(w\in\mathbf R\) and every
\(K\in\DW(X_0,E)\),
\[
 \boxed{\qquad
       \IntWtGe(K,w)\quad\Longrightarrow\quad
       \IntWtGe(Rf_{*,\mathrm W}K,w).
 \qquad}
\]
\dep{Verdier exchange, the definition of lower weights, and Deligne's
upper-weight theorem applied to \(D_XK\).}
\end{theorem}

\begin{corollary}[Ordinary lower bound for integral inputs]
Under the hypotheses of the lower-weight theorem, for every
\(w\in\mathbf R\) and \(K\in\DW(X_0,E)\),
\[
          \IntWtGe(K,w)\quad\Longrightarrow\quad
          \WtGe(Rf_{*,\mathrm W}K,w).
\]
\end{corollary}

\begin{theorem}[Smooth lisse lower-weight bridge]
Assume the finite-level exceptional-duality hypotheses.  Let \(X_0\) be
smooth of pure dimension \(d\) and separated over \(k\), and let
\(\mathcal L\) be a lisse Weil sheaf punctually \(\iota\)-pure of an
integral weight \(w\in\mathbf Z\).  Smooth Poincare duality gives a strict
Frobenius-equivariant isomorphism
\[
                 D_X\mathcal L\xrightarrow{\sim}
                 \mathcal L^\vee(d)[2d].
\]
Consequently
\[
             \IntWtLe(\mathcal L,w),\qquad
             \IntWtGe(\mathcal L,w),
\]
so \(\mathcal L\), regarded as a complex in degree zero, is pure of weight
\(w\) in the upper/lower sense used in this note.  The same assertion holds
after finite extension of the coefficient field or the finite ground
field.
\dep{The smooth identification \(a_X^!E\simeq E(d)[2d]\), strict
Frobenius compatibility of exceptional pullback, the dual and Tate weight
formulas, and the definitions of \(\IntWtLe\) and \(\IntWtGe\).}
\end{theorem}

\begin{corollary}[Proper direct image preserves purity]
Assume the finite-level exceptional-duality hypotheses.  Let
\(f_0:X_0\to Y_0\) be proper with \(X_0\) and \(Y_0\) separated over
\(k\), let \(w\in\mathbf Z\), and let \(K\in\DW(X_0,E)\) satisfy both
\(\IntWtLe(K,w)\) and \(\IntWtGe(K,w)\).  The Frobenius-equivariant
comparison \(Rf_{!,\mathrm W}\xrightarrow{\sim}Rf_{*,\mathrm W}\) and the
upper and lower direct-image theorems show that
\(Rf_{*,\mathrm W}K\) satisfies the same two conditions and is pure of weight
\(w\) in the upper/lower BBD sense.  If
\(Y_0=\Spec k\), this is equivalent to
\[
                 H^i(Rf_*K)\text{ being pure of weight }w+i
\]
for every \(i\), because Verdier duality on a point is ordinary linear
duality.  No such equivalence with punctual purity of every standard
cohomology sheaf is asserted for a general target.
\dep{The upper theorem, the lower theorem, and proper comparison.}
\end{corollary}

\begin{corollary}[Smooth proper cohomology]
Assume the finite-level exceptional-duality hypotheses.  Let
\(f_0:X_0\to Y_0\) be smooth and proper, with \(X_0\) and \(Y_0\)
separated over \(k\), and let \(\mathcal L\) be a
lisse Weil sheaf punctually \(\iota\)-pure of an integral weight
\(w\in\mathbf Z\).  Then, for every \(i\in\mathbf Z\),
\(R^if_*\mathcal L\) is lisse and punctually \(\iota\)-pure of weight
\(w+i\).  At a closed point
\(y\in|Y_0|\), proper smooth base change identifies its Frobenius module with
\[
 H^i(X_{\bar y},\mathcal L|_{X_{\bar y}}),
\]
which is pure of weight \(w+i\).  In particular, if \(X_0\) is smooth and
proper over \(k\),
\[
                    H^i(X_{\bar k},E)
       \text{ is pure of weight }i.
\]
\dep{The smooth lisse lower-weight bridge on every closed geometric fiber,
the point-valued case of proper direct-image purity, smooth proper base
change, lissity of \(R^if_*\mathcal L\), and closed-point detection for a
lisse Weil sheaf.}
\end{corollary}

\section{References}

\begin{itemize}
\item P. Deligne, \emph{La conjecture de Weil II}, Publ. Math. IHES
      \textbf{52} (1980), especially Theorems 3.2.3 and 3.3.1 and
      Corollaries 3.3.4--3.3.6.
\item G. Laumon, \emph{Transformation de Fourier, constantes d'equations
      fonctionnelles et conjecture de Weil}, Publ. Math. IHES
      \textbf{65} (1987), Sections 4.1--4.4.
\item Stacks Project, chapters on constructible etale sheaves, proper base
      change, compactly supported cohomology, normalization, Zariski's Main
      Theorem, and limits of schemes.
\end{itemize}

\end{document}
